Explain the general challenge for which the classical linear
model provides an example


Key goal:
The key goal is to understand a relationship between a response variable \( y \) of primary interest and a set of regressors/covariates or explanatory variables 
\( x_1, x_2, \ldots, x_k \), i.e.,
\[
 y = f(x_1, x_2, \ldots, x_k)
\]
with some unknown deterministic function \( f \)
In practice, it is often reasonable that the relationship involves some error due to some noise.

Specific classes of models:

Specific models (or classes of models) impose restriction on the class of admissible functions, the way errors enter, and the structure of the error.
Within classes of models, due to their particular properties, statistical tools can often be developed.
A key initial question is, of course, if a particular model is appropriate and how we can test this.
Classical Linear Models are one particular class of models.

What is the special case of the linear model?

Classical Linear Model

\[ y = f(x_1, x_2, \ldots, x_k) + \epsilon \]

Systematic component:
- The unknown function \( f \) is called systematic component.
- This is assumed to be a linear combination of the covariates:

\[ f(x_1, x_2, \ldots, x_k) = \beta_0 + \beta_1 x_1 + \cdots + \beta_k x_k \]

for unknown parameters \( \beta_0, \beta_1, \ldots, \beta_k \).

The model is conveniently expressed in vector notation:

\[ f(\mathbf{x}) = \mathbf{x}^\top \boldsymbol{\beta}, \quad \mathbf{x} = (1, x_1, \ldots, x_k)^\top, \quad \boldsymbol{\beta} = (\beta_0, \ldots, \beta_k)^\top. \]

Additive errors:
In the linear model it is assumed that errors are additive:

\[ y = \mathbf{x}^\top \boldsymbol{\beta} + \epsilon \]



For a fixed number of covariate vectors, state the linear model in vector notation
(include terms like design matrix, normal regression)


In order to determine the unknown parameters \(\beta\), we collect data
\[ y_i, \; \mathbf{x}_i = (1, x_{i1}, \ldots, x_{ik})^\top, \]
with error observations according to
\[ y_i = \mathbf{x}_i^\top \boldsymbol{\beta} + \epsilon_i, \quad i = 1, 2, \ldots, n. \]
Vector notation is appropriate:

\[ 
\mathbf{y} = \begin{pmatrix}
y_1 \\
\vdots \\
y_n
\end{pmatrix}, \quad \boldsymbol{\epsilon} = \begin{pmatrix}
\epsilon_1 \\
\vdots \\
\epsilon_n
\end{pmatrix}, \quad \mathbf{X} = \begin{pmatrix}
\mathbf{x}_1^\top \\
\vdots \\
\mathbf{x}_n^\top
\end{pmatrix} = \begin{pmatrix}
1 & x_{11} & \cdots & x_{1k} \\
\vdots & \vdots & \ddots & \vdots \\
1 & x_{n1} & \cdots & x_{nk}
\end{pmatrix}.
\]

\(\mathbf{X}\) is the design matrix. One assumes that \(\text{rk} \mathbf{X} = k + 1 = p \leq n\).


The model
\[ y = X \boldsymbol{\beta} + \epsilon \]

is called the classical linear regression model, if the following assumptions hold:

\[ \mathbb{E}(\epsilon) = 0 \]
\[ \text{Cov}(\epsilon) = \sigma^2 I \]
this correlation structure is also referred to as homoscedastic
while different variances would be heteroscedastic.
One says that these assumptions about \( \epsilon \) are conditionally on \( X \).
\[ \text{rk}(X) = k + 1 = p \leq n \]

If additionally errors are Gaussian, the model is called the classical normal regression model.
For stochastic covariates all properties are understood conditionally on the design matrix \(X\).


Discuss the estimation problem and issues related to misspecification.


Linearity of Covariate Effects
Linear models are linear in their parameters, but covariates can be modified non-linearly.

An example is 
\[ y_i = \beta_0 + \beta_1 \log(z_i) + \epsilon_i, \quad x_i = \log(z_i). \]

Homoscedastic Error Variances
Residuals can be used to detect heteroscedasticity.
If this is ignored, variances of \(\hat{\beta}\) are incorrectly estimated, causing trouble with hypothesis tests and confidence intervals.

Uncorrelated Errors
Autocorrelation can also be detected by inspecting the residuals.
It indicates that the model is incorrectly specified, e.g., covariates are missing or should enter nonlinearly. 
E.g. we include \( x_i \) linearly but the relation is actually \( y_i = \sin(x_i) + \epsilon_i \).
or \( y_i = x_{1_{i}} + x_{2_{i}} + c + \epsilon_i \), but we only specify \( y_i = x_{1_{i}} + \epsilon_i \).
The statistical analysis should be adjusted in the presence of autocorrelation.

Additivity of Errors
Many different ways can be imagined how errors enter the model.
Most common, are additive and multiplicative errors. Multiplicative errors may be transformed into additive errors by taking logarithms:

\[ y_i = \exp(\boldsymbol{\beta}^\top \mathbf{x}_i + \epsilon_i) \quad \Longleftrightarrow \quad \log y_i = \boldsymbol{\beta}^\top \mathbf{x}_i + \epsilon_i \]

iNCLUDE IMAGES FROM PAGE 104-109


How can categorial variables/Continuous Covariates be captured?

If a continuous covariate has approximately a nonlinear effect \(\beta_1 f(z)\) with known function \(f\), then the model
\[ y_i = \beta_0 + \beta_1 f(z_i) + \cdots + \epsilon_i \]
can be replaced by the linear regression model
\[ y_i = \beta_0 + \beta_1 x_i + \cdots + \epsilon_i, \quad x_i = f(z_i). \]
If covariates of the new models should be centered around zero, alternatively one may set
\[ x_i = f(z_i) - \bar{f}, \quad \bar{f} = \frac{1}{n} \sum_{i=1}^n f(z_i) \]
with new \(\beta_0, \beta_1, \ldots,\) and the multiple of \(\bar{f}\) is absorbed into the new \(\beta_0\).
E.g. the covariate enters polynomially as \( \beta_1 z + \beta_2 z^2 + \dots \)
then we would transform this into a linear form by setting \( x_{i,j} = z^j \) e.g.
\( y_i = \beta_0 + \beta_1 x_{i,1} + \beta_2 x_{i,2} \dots \)


Dummy Coding for Categorial Covariates

Suppose that \( x_i \) are observations of a categorial covariate with categories \(\{1, 2, \ldots, c\}\). Dummy coding sets

\[
x_{ik} =
\begin{cases} 
1 & x_i = k \\
0 & \text{else} 
\end{cases} \quad k = 1, 2, \ldots, c - 1, \quad i = 1, 2, \ldots, n.
\]

This means that \( x_{ik} = 1 \), iff \( x_i \) is of category \( k \), and otherwise zero. Dummies are indicators that display the category. 
Category \( c \) is omitted, since this \( c \) can already be identified by verifying that all dummies are zero. 
Dummies allow to include the categorial variable into the regression model:

\[
y_i = \beta_0 + \beta_{i1} x_{i1} + \cdots + \beta_{i, c-1} x_{i, c-1} + \cdots + \epsilon_i, \quad i = 1, 2, \ldots, n.
\]

Remark: Other codings for categorial variables, such as effect coding, are also sometimes used. Effect coding refers to

\[
x_{ik} =
\begin{cases} 
1 & x_i = k \\
-1 & x_i = c \\
0 & \text{else} 
\end{cases} \quad k = 1, 2, \ldots, c - 1, \quad i = 1, 2, \ldots, n.
\]


What is the method of least squares? What is median regression? Derive the LSE
of β?

The classical method to estimate the unknown regression coefficients \(\beta\) consists in minimizing the sum of the squared deviations

\[
LS(\beta) = \sum_{i=1}^{n} (y_i - \mathbf{x}_i^\top \beta)^2 = \mathbf{\varepsilon}^\top \mathbf{\varepsilon}
\]

with respect to \(\beta \in \mathbb{R}^p\).

The solution is the same as the maximum likelihood estimator with Gaussian error terms.

The estimator, commonly called least squares estimator (LSE), has the disadvantage to depend heavily on outliers. 
But it is easy to compute, and its statistical properties have long been understood.

An alternative is, e.g., median regression, a special case of quantile regression, that minimizes with respect to \(\beta \in \mathbb{R}^p\) the sum of absolute deviations

\[
SM(\beta) = \sum_{i=1}^{n} |y_i - \mathbf{x}_i^\top \beta| = \sum_{i=1}^{n} |\varepsilon_i|.
\]


Show how LSE is related to MLE in the Gaussian case.

Minimization of \(LS\) can be achieved by inspecting the first-order conditions:
\[
\frac{\partial LS(\beta)}{\partial \beta} \bigg|_{\beta = \hat{\beta}} = -2\mathbf{X}^\top \mathbf{y} + 2 \mathbf{X}^\top \mathbf{X} \hat{\beta} \stackrel{!}{=} 0.
\]

This is equivalent to the normal equations:
\[
\mathbf{X}^\top \mathbf{X} \hat{\beta} = \mathbf{X}^\top \mathbf{y}.
\]

Due to the condition of full rank, the normal equations possess a unique solution, the least squares estimator,
\[
\hat{\beta} = (\mathbf{X}^\top \mathbf{X})^{-1} \mathbf{X}^\top \mathbf{y}.
\]


Under the condition of normally distributed errors, i.e., \(\varepsilon \sim N(0, \sigma^2 I)\), 
the maximum likelihood estimator (MLE) equals the LSE.
This can be verified by minimizing the likelihood

\[
L(\beta, \sigma^2) = \frac{1}{(2\pi\sigma^2)^{n/2}} \exp \left( -\frac{1}{2\sigma^2} (y - \mathbf{X}\beta)^\top (y - \mathbf{X}\beta) \right),
\]

or, equivalently, the log-likelihood

\[
l(\beta, \sigma^2) = (\text{terms independent of } \beta) - \frac{1}{2\sigma^2} (y - \mathbf{X}\beta)^\top (y - \mathbf{X}\beta).
\]


How is the prediction/hat matrix \( H \) defined? State its properties and relation to residuals and mean.


The estimated (conditional) mean of \(y\) is

\[
\hat{y} = \mathbb{E}(\hat{y}) = \mathbf{X} \hat{\beta} = \mathbf{H}y, \quad \mathbf{H} = \mathbf{X} (\mathbf{X}^\top \mathbf{X})^{-1} \mathbf{X}^\top.
\]

The matrix \(\mathbf{H}\) is called the prediction matrix or hat matrix.
The residuals can be computed as \(\hat{\varepsilon} = (\mathbf{I} - \mathbf{H})y\).
Properties of the prediction matrix:
The prediction matrix \(\mathbf{H} = \mathbf{X} (\mathbf{X}^\top \mathbf{X})^{-1} \mathbf{X}^\top\) is symmetric and idempotent with \(\text{rk}(\mathbf{H}) = p\).

The matrix \(\mathbf{I} - \mathbf{H}\) is symmetric and idempotent, with \(\text{rk}(\mathbf{I} - \mathbf{H}) = n - p\).

Since \( \hat{\mathbf{y}} \in \text{im}(\mathbf{X}), \hat{\mathbf{\epsilon}} \bot \text{im}(\mathbf{X}) \implies \hat{\mathbf{y}} \bot \hat{\mathbf{\epsilon}} \)
it follows \( \text{proj}_{\text{im}(\mathbf{X})}(\mathbf{y})= \hat{\mathbf{y}} \). In particular \( \text{proj}_{\text{im}(\mathbf{X})} = H \)

\( \hat{\mathbf{\epsilon}} = \mathbf{y} - \hat{\mathbf{y}} = y - \mathbf{X}\hat{\beta} = (\mathbf{I} - \mathbf{X} (\mathbf{X}^\top \mathbf{X})^{-1} \mathbf{X}^\top\mathbf{X})\mathbf{y} = (\mathbf{I} - \mathbf{H})\mathbf{y} \).


Which variance estimator is usually used? How does it differ from the MLE obtained in the Gaussian case?

A heuristics for finding an estimator of the error variance can be constructed in the setting with Gaussian errors.
Maximum Likelihood

From the log-likelihood we can find an ML-estimator

\[
\hat{\sigma}^2_{ML} = \frac{\hat{\varepsilon}^\top \hat{\varepsilon}}{n}
\]

But this estimator is biased, i.e., \(\mathbb{E}(\hat{\sigma}^2_{ML}) = \frac{n-p}{n} \cdot \sigma^2\).
This suggests the commonly used estimator

\[
\hat{\sigma}^2 = \frac{1}{n-p} \hat{\varepsilon}^\top \hat{\varepsilon}
\]

Restricted Maximum Likelihood
In the Gaussian framework, \(\hat{\sigma}^2\) is a restricted maximum likelihood estimator (RMLE) that minimizes the marginal likelihood


What are expectation and variance of \( \hat{β} \)? 
State the corresponding Gauss-Markov Theorem. What do we obtain in the Gaussian case?

\[
L(\sigma^2) = \int L(\beta, \sigma^2) d\beta
\]

The following properties hold if the model is correctly specified, but without any specific distributional assumptions:

\(\mathbb{E}(\hat{\beta}) = \beta\) (unbiased)

\(
\text{Cov}(\hat{\beta}) = \sigma^2 (\mathbf{X}^\top \mathbf{X})^{-1} \quad \text{with estimator} \quad \widehat{\text{Cov}}(\hat{\beta}) = \hat{\sigma}^2 (\mathbf{X}^\top \mathbf{X})^{-1}, \quad \hat{\sigma}^2 = \frac{1}{n-p} \hat{\varepsilon}^\top \hat{\varepsilon}
\)

Gauss-Markov-Theorem: The LSE has the smallest variance among all linear and unbiased estimators of \(\beta\).

Suppose that
\[
  \lim_{n \to \infty} \frac{1}{n} \mathbf{X}_n^T \mathbf{X}_n = \mathbf{V} \qquad (1)
\]
with \( \mathbf{V} \in \mathbb{R}^{p,p} \) positive definite with \( p = k + 1 \) and \( \mathbf{X}_n \) is the design matrix with
\( n \)-observations.
then it holds:
The LS estimator \( \hat{\beta}_n \) and REML estimator \( \hat{\sigma}^2_n \) are consistent for \( \beta \) and \( \sigma^2 \), respectively.
The LS estimator \( \hat{\beta}_n \) is asymptotically normal with
\[
\sqrt{n} (\hat{\beta}_n - \beta) \xrightarrow{d} N(0, \sigma^2 \mathbf{V}^{-1}).
\]
This implies that for large \( n \) the estimator \( \hat{\beta}_n \) possesses approximately the normal distribution
\[
  N\left( \beta, \hat{\sigma}^2_n (\mathbf{X}_n^T \mathbf{X}_n)^{-1} \right).
\]

The assumption (1) is satisfied if the covariates \(\mathbf{x}_i\) are i.i.d.
A counterexamples are covariates with a trend e.g. in 1D, \( x_i = i \implies \frac{1}{n}\frac{1}{n} X_n^T X_n = \sum_{i=1}^n i^2 \), 
The weaker assumptions,
\( (\mathbf{X}_n^\top\mathbf{X}_n)^{-1} \xrightarrow 0 \) gives consistency
\( \sup_i \mathbf{x}_i^\top(\mathbf{X}_n^\top\mathbf{X}_n)^{-1}\mathbf{x}_i \xrightarrow 0 \) gives asymptotic normality
actually give the same result and resolve this issue.

The following properties hold if additionally \(\varepsilon \sim N(0, \sigma^2 \mathbf{I})\):
Distribution of the response variable: \(y \sim N(\mathbf{X}\beta, \sigma^2 \mathbf{I})\)
Distribution of LSE: \(\hat{\beta} \sim N(\beta, \sigma^2 (\mathbf{X}^\top \mathbf{X})^{-1})\)
Distribution of weighted distance:

\[
\frac{(\hat{\beta} - \beta)^\top (\mathbf{X}^\top \mathbf{X})(\hat{\beta} - \beta)}{\sigma^2} \sim \chi^2_p
\]

Proof of unbiasedness

\((\mathbf{X}^\top \mathbf{X})^{-1}\mathbf{X}^\top\) is constant 

\[
\mathbb{E} \left[ \hat{\beta} \right] = \mathbb{E} \left[ (\mathbf{X}^\top \mathbf{X})^{-1}\mathbf{X}^\top \mathbf{y} \right] = (\mathbf{X}^\top \mathbf{X})^{-1}\mathbf{X}^\top \mathbb{E}[\mathbf{y}].
\]

Substituting in \(\mathbf{X}\beta + \epsilon \text{ for } \mathbf{y}\)

\[
\mathbb{E} \left[ \hat{\beta} \right] = (\mathbf{X}^\top \mathbf{X})^{-1}\mathbf{X}^\top \mathbb{E}[\mathbf{X}\beta + \epsilon]
= (\mathbf{X}^\top \mathbf{X})^{-1}\mathbf{X}^\top (\mathbf{X}\beta + \mathbb{E}[\epsilon])
= (\mathbf{X}^\top \mathbf{X})^{-1}\mathbf{X}^\top (\mathbf{X}\beta + 0)
= I_p \beta
= \beta.
\]

Derivation of variance:

Since for a constant matrix \( \mathbf{A} \), we have that \( \text{Var}(\mathbf{AY}) = \mathbf{A} \text{Var}(\mathbf{Y}) \mathbf{A}^\top \), 
it follows that

\[
\text{Var}(\hat{\beta}) = \text{Var} \left( (\mathbf{X}^\top \mathbf{X})^{-1} \mathbf{X}^\top \mathbf{y} \right)
= (\mathbf{X}^\top \mathbf{X})^{-1} \mathbf{X}^\top \text{Var}(\mathbf{y}) \left( (\mathbf{X}^\top \mathbf{X})^{-1} \mathbf{X}^\top \right)^\top
\]

Substituting in \( \mathbf{X}\beta + \epsilon \text{ for } \mathbf{y} \) and noting that \( \mathbf{X}\beta \) is a constant, we have that

\[
\text{Var}(\hat{\beta}) = (\mathbf{X}^\top \mathbf{X})^{-1} \mathbf{X}^\top \text{Var}(\mathbf{X}\beta + \epsilon) \mathbf{X} (\mathbf{X}^\top \mathbf{X})^{-1}
= (\mathbf{X}^\top \mathbf{X})^{-1} \mathbf{X}^\top \text{Var}(\epsilon) \mathbf{X} (\mathbf{X}^\top \mathbf{X})^{-1}
\]
\[
= (\mathbf{X}^\top \mathbf{X})^{-1} \mathbf{X}^\top \sigma^2 \mathbf{I}_n \mathbf{X} (\mathbf{X}^\top \mathbf{X})^{-1}
= \sigma^2 (\mathbf{X}^\top \mathbf{X})^{-1} \mathbf{X}^\top \mathbf{X} (\mathbf{X}^\top \mathbf{X})^{-1}
= \sigma^2 (\mathbf{X}^\top \mathbf{X})^{-1}.
\]


Are there limit theorems available in the general case that admit an approximation?
Can you state details?
